\textbf{Relación con la suma y la resta} (\(a, b \ge 0\)).
\begin{props}
  \item \(a + b = (a \bxor b) + 2\,(a \band b)\): el XOR es la suma sin acarreo y \(a \band b\) marca dónde hay acarreo.
  \item \(a + b = (a \bor b) + (a \band b)\)
  \item \(a \bxor b = (a \bor b) - (a \band b)\)
  \item \(a \bor b = (a \bxor b) + (a \band b) = (a \bxor b) \bor (a \band b)\)
  \item \(a \band b = (a \bor b) \bxor (a \bxor b)\)
  \item \(a \bxor b \le a + b\), con igualdad \(\iff a \band b = 0\).
  \item \(\max(a,b) \le a \bor b \le a + b\), \(\quad a \band b \le \min(a,b)\).
\end{props}

\textbf{Submáscaras y conteo de bits} (\(\operatorname{pc}\) = cantidad de bits en 1).
\begin{props}
  \item \(a\) es submáscara de \(b\) (todo bit encendido de \(a\) está encendido en \(b\)) \(\iff a \band b = a \iff a \bor b = b\).
  \item \(\operatorname{pc}(a \bor b) = \operatorname{pc}(a) + \operatorname{pc}(b) - \operatorname{pc}(a \band b)\)
  \item \(\operatorname{pc}(a \bxor b) = \operatorname{pc}(a) + \operatorname{pc}(b) - 2\operatorname{pc}(a \band b)\)
  \item La distancia de Hamming entre \(a\) y \(b\) (bits distintos) es \(\operatorname{pc}(a \bxor b)\).
\end{props}

\textbf{Desplazamientos (shifts).} Notación: \(\ll\) shift izquierda, \(\gg\) shift derecha.
\begin{props}
  \item \(a \ll k = a \cdot 2^k\), \(\quad a \gg k = \lfloor a / 2^k \rfloor\). El piso también aplica a negativos: \(-5 \gg 1 = -3\).
  \item \((a \ll j) \ll k = a \ll (j+k)\), \(\quad (a \gg j) \gg k = a \gg (j+k)\).
  \item Los shifts distribuyen sobre AND, OR y XOR: \((a \band b) \ll k = (a \ll k) \band (b \ll k)\), y lo mismo para \(\bor\), \(\bxor\) y \(\gg\).
  \item \(2^k = 1 \ll k\) y \(2^k - 1 = (1 \ll k) - 1\) (máscara de \(k\) bits en 1).
  \item \(a \bmod 2^k = a \band ((1 \ll k) - 1)\).
  \item \(a\) es par \(\iff a \band 1 = 0\). \(a\) es múltiplo de \(2^k \iff a \band ((1 \ll k) - 1) = 0\).
\end{props}

\textbf{Precedencia de operadores en Python} (de mayor a menor): \(\bnot\) (unario), \(*\ /\ \%\), \(+\ -\), \(\ll\ \gg\), \(\band\), \(\bxor\), \(\bor\), comparaciones (\(==\), \(<\), \dots). Dos consecuencias frecuentes:
\begin{props}
  \item Los shifts tienen \textbf{menor} precedencia que \(+\) y \(-\): \texttt{1 <{}< k + 1} es \texttt{1 <{}< (k + 1)}, no \texttt{(1 <{}< k) + 1}.
  \item A diferencia de C/C++, en Python \texttt{a \& b == c} se evalúa como \texttt{(a \& b) == c}. Aun así conviene poner paréntesis para no depender de ello.
\end{props}
